\subsection{3. Mục tiêu đề tài}

\indent\indent Mục tiêu tổng quát của đề tài là nghiên cứu và xây dựng hệ thống đánh giá rủi ro và phân loại tài khoản Sybil trên mạng xã hội phi tập trung, ứng dụng mô hình đồ thị trí tuệ nhận thức ràng buộc. Đề tài hướng tới việc giải quyết bài toán thiếu hụt nhãn dữ liệu thực tế, đồng thời nâng cao hiệu quả phân lớp đa thông qua việc học biểu diễn các mối quan hệ phức tạp từ hành vi tương tác, nội dung ngữ nghĩa và các ràng buộc đồ thị, đặc biệt và các quan hệ mang tính sở hữu trên chuỗi khối. Các mục tiêu cụ thể mà đề tài hướng đến bao gồm:

\begin{itemize}[label={--}, itemsep=1pt]
    \item \textbf{Xây dựng tập dữ liệu đa đặc trưng:} Trích xuất và tiền xử lý dữ liệu thực tế từ Lens Protocol qua các mốc thời gian đa dạng. Khai phá các đặc trưng thống kê hành vi và sử dụng mô hình ngôn ngữ (Sentence-BERT) để nhúng ngữ nghĩa hồ sơ cá nhân, thiết lập không gian đặc trưng tiêu chuẩn phục vụ huấn luyện mô hình học sâu.
    
    \item \textbf{Mô hình hóa đồ thị đa tầng nhận thức ràng buộc (CMG):} Xây dựng cấu trúc đồ thị đa quan hệ với 4 tầng kiến trúc chuyên biệt (Theo dõi, Tương tác, Tương đồng và Đồng sở hữu). Tích hợp cơ chế tinh chỉnh trọng số cơ sở nhằm định hướng mô hình nắm bắt chính xác hơn các ràng buộc về tính sở hữu và cấu trúc của các cụm Sybil.
    
    \item \textbf{Đề xuất quy trình gán nhãn tự động:} Khắc phục sự thiếu hụt dữ liệu gán nhãn. Thông qua quy trình ứng dụng GAE và thuật toán phân cụm kết hợp hệ thống luật để tự động phân loại và gán nhãn giả lập theo 4 mức độ rủi ro (\texttt{BENIGN, LOW\_RISK, HIGH\_RISK, MALICIOUS}).
    
    \item \textbf{Nghiên cứu đối sánh các kiến trúc GNN:} Triển khai và đánh giá hiệu năng của các họ mô hình GNN tiên tiến để kiểm chứng khả năng học biểu diễn và phân lớp đa nhãn trực tiếp trên cấu trúc đồ thị của nghiên cứu này.
    
    \item \textbf{Tối ưu hiệu năng với kiến trúc Decoupled GNN-ML:} Đề xuất và khảo sát các tổ hợp mô hình tách rời, sử dụng GNN như bộ trích xuất đặc trưng kết hợp cùng các thuật toán học máy truyền thống. Việc này nhằm thiết lập ranh giới phân loại rõ nét hơn giữa 4 nhóm rủi ro, từ đó tối ưu hóa độ chính xác.
    
    \item \textbf{Phát triển hệ thống thử nghiệm có tính diễn giải:} Đóng gói các thuật toán thành một ứng dụng web hoàn chỉnh phục vụ việc tra cứu mức độ rủi ro thực tế. Tích hợp khả năng trực quan hóa đồ thị cho mục đích dễ dàng giải thích, nhằm giải quyết bài toán "hộp đen" của AI và đảm bảo tính minh bạch.
\end{itemize}